% !TEX root = ../discretemath.tex

\section{Моменты случайной величины}

\subsection{Основные определения}


\begin{Definition}{Моменты случайной величины}{}
	Пусть $\xi$ — случайная величина, $k \in \mathbb N$.
	
	\begin{itemize}
		\item \textbf{Начальный момент} порядка $k$:
		\[
		\mu'_k = \mathbb{E}[\xi^k].
		\]
		При $k=1$ начальный момент $\mu'_1 = \mathbb{E}[\xi]$
		называется математическим ожиданием.
		
		\item \textbf{Центральный момент} порядка $k \ge 2$:
		\[
		\mu_k = \mathbb{E}\big[(\xi - \mathbb{E}[\xi])^k\big].
		\]
		По определению $\mu_1 = 0$, а величина $\mu_2 = \mathbb{E}\big[(\xi - \mathbb{E}[\xi])^2\big]$ называется дисперсией.
	\end{itemize}
\end{Definition}

\subsection{Связь центральных и начальных моментов}

\begin{Theorem}{Выражение центрального момента через начальные}{central-via-initial}
	Для любого $k \geq 2$
	\[
	\mu_k = \sum_{s=0}^{k} (-1)^s \binom{k}{s} \bigl(\mathbb{E}[\xi]\bigr)^s \cdot \mu'_{k-s}.
	\]
\end{Theorem}

\begin{proof}
Раскроем определение центрального момента с помощью формулы бинома Ньютона:
\begin{align*}
	\mu_k
	&= \mathbb{E}\!\left[(\xi - \mathbb{E}[\xi])^k\right]
	= \mathbb{E}\!\left[\sum_{s=0}^{k} (-1)^s \binom{k}{s} (\mathbb{E}[\xi])^s \cdot \xi^{k-s}\right] \\
	&= \sum_{s=0}^{k} (-1)^s \binom{k}{s} (\mathbb{E}[\xi])^s \cdot \mathbb{E}[\xi^{k-s}]
	= \sum_{s=0}^{k} (-1)^s \binom{k}{s} (\mathbb{E}[\xi])^s \cdot \mu'_{k-s},
\end{align*}
где на последнем шаге использована линейность математического ожидания и
тот факт, что $\mathbb{E}[\xi^{k-s}] = \mu'_{k-s}$.
\end{proof}

\begin{Corollary}{}{}
	Поскольку $\mathbb{E}[\xi] = \mu'_1$, формула принимает вид
	\[
	\mu_k = \sum_{s=0}^{k} (-1)^s \binom{k}{s} (\mu'_1)^s \cdot \mu'_{k-s}.
	\]
\end{Corollary}

\begin{Corollary}{}{}
	Каждый центральный момент $\mu_k$ выражается через начальные моменты $\mu'_1, \ldots, \mu'_k$,
	и наоборот: каждый начальный момент $\mu'_k$ выражается через $\mu'_1$ и центральные моменты $\mu_2, \ldots, \mu_k$.
\end{Corollary}

\begin{Remark}{}{}
	Это возможно потому, что наборы
	$\{1,\xi,\xi^2,\ldots,\xi^k\}$ и
	$\{1,(\xi-\mathbb E\xi),(\xi-\mathbb E\xi)^2,\ldots,(\xi-\mathbb E\xi)^k\}$
	являются базисами одной и той же линейной оболочки.
\end{Remark}

\subsection{Единственность распределения по моментам}

\begin{Theorem}{}{}
	Пусть $\xi$ и $\beta$ — две случайные величины над конечным вероятностным пространством $P$,
	и все их моменты совпадают: $\mu'_k(\xi) = \mu'_k(\beta)$ для всех $k$.
	Тогда $\xi \overset{d}{=} \beta$ (случайные величины имеют одинаковое распределение).
\end{Theorem}

\begin{proof}
Пусть $x_1, \ldots, x_n$ — все возможные значения $\xi$ и $\beta$ (объединённое множество значений).
Введём обозначения:
\[
d_\xi(i) = P[\xi = x_i], \qquad \sum_{i=1}^n d_\xi(i) = 1,
\]
и аналогично $d_\beta(i) = P[\beta = x_i]$.

Рассмотрим матрицу Вандермонда:
\[
M =
\begin{pmatrix}
	1      & 1      & \cdots & 1      \\
	x_1    & x_2    & \cdots & x_n    \\
	x_1^2  & x_2^2  & \cdots & x_n^2  \\
	\vdots & \vdots & \ddots & \vdots \\
	x_1^{n-1} & x_2^{n-1} & \cdots & x_n^{n-1}
\end{pmatrix}.
\]
Тогда:
\[
M \cdot
\begin{pmatrix} d_\xi(1) \\ \vdots \\ d_\xi(n) \end{pmatrix}
=
\begin{pmatrix}
	\sum d_\xi(i) \\
	\sum x_i\, d_\xi(i) \\
	\vdots \\
	\sum x_i^{n-1} d_\xi(i)
\end{pmatrix}
=
\begin{pmatrix}
	1 \\ \mu'_1(\xi) \\ \vdots \\ \mu'_{n-1}(\xi)
\end{pmatrix}.
\]
Поскольку моменты $\xi$ и $\beta$ совпадают, имеем
\[
M \cdot \begin{pmatrix} d_\xi(1) \\ \vdots \\ d_\xi(n) \end{pmatrix}
=
M \cdot \begin{pmatrix} d_\beta(1) \\ \vdots \\ d_\beta(n) \end{pmatrix},
\]
откуда
\[
M \cdot
\begin{pmatrix} d_\xi(1) - d_\beta(1) \\ \vdots \\ d_\xi(n) - d_\beta(n) \end{pmatrix}
= \mathbf{0}.
\]
Покажем, что $\det M \neq 0$. Определитель $M$ есть однородный многочлен от $x_1, \ldots, x_n$
степени $\frac{n(n-1)}{2}$. Поскольку $\det M$ обращается в нуль при $x_i = x_j$ ($i \neq j$),
каждый множитель $(x_i - x_j)$ делит $\det M$. Так как
\[
\deg(\det M) = \deg\!\prod_{i < j}(x_i - x_j) = \frac{n(n-1)}{2},
\]
то $\det M = C \cdot \prod_{i < j}(x_i - x_j)$ для некоторой константы $C$.

Для нахождения $C$ заметим, что при делении $\det M$ последовательно на множители
$(x_i - x_j)$ получается константа, и коэффициент при мономе
$x_1^0 x_2^1 \cdots x_n^{n-1}$ в $\det M$ равен $1$
(он порождается единственной перестановкой в определителе).
Аналогично, коэффициент при том же мономе в $\prod_{i<j}(x_i - x_j)$ равен $1$,
откуда $C = 1$ и
\[
\det M = \prod_{i < j}(x_i - x_j) \neq 0
\]
при попарно различных $x_1, \ldots, x_n$.

Следовательно, матрица $M$ невырождена, и из $M \cdot (d_\xi - d_\beta) = \mathbf{0}$
получаем $d_\xi(i) = d_\beta(i)$ для всех $i$, то есть $\xi$ и $\beta$
имеют одно и то же распределение. 
\end{proof}
